% Generated draft: requires mathematical and source review.
\documentclass[11pt]{article}
\usepackage[margin=25mm]{geometry}
\usepackage{fontspec}
\setmainfont{Latin Modern Roman}
\usepackage{amsmath,amssymb,mathtools,unicode-math}
\setmathfont{Latin Modern Math}
\usepackage{enumitem,needspace,xurl,hyperref}
\setlength{\parindent}{0pt}
\setlength{\parskip}{5pt}
\setlength{\emergencystretch}{3em}
\newcommand{\N}{\mathbb{N}}
\newcommand{\Z}{\mathbb{Z}}
\newcommand{\Q}{\mathbb{Q}}
\newcommand{\R}{\mathbb{R}}
\newcommand{\C}{\mathbb{C}}
\newcommand{\F}{\mathbb{F}}
\newcommand{\A}{\mathbb{A}}
\newcommand{\PP}{\mathbb{P}}
\newcommand{\E}{\mathbb{E}}
\newcommand{\eps}{\varepsilon}
\newcommand{\dd}{\,\mathrm{d}}
\begin{document}
\section*{UnsolvedMath ID: 5300070}
\textbf{Source problem code:} catalogue entry
\subsection{Definitions and mathematical statement}
In the usual complex-polynomial setting, let \(f\in\mathbb C[z]\) have degree \(d\), and for a real parameter \(0<h\le1\) define the relaxed Newton rational map
\[
 N_{h,f}(z)=z-h\frac{f(z)}{f'(z)},
\]
with common factors cancelled and the map extended to the Riemann sphere where appropriate. The ordinary Newton map is \(N_{1,f}\). A periodic orbit is attracting when its multiplier, computed in a local coordinate, has absolute value less than \(1\). Call \(f\) bad for parameter \(h\) if \(N_{h,f}\) has an attracting periodic orbit that is not a root fixed point of \(f\).

The source asks that every bad degree-\(d\) polynomial \(f_1\) occur as the endpoint of a one-parameter family \((f_h)_{0<h\le1}\), with that prescribed value at \(h=1\), such that
\[
 \forall h\in(0,1]:\quad
 N_{h,f_h}\text{ has an attracting periodic orbit other than a root},
\]
and, as \(h\downarrow0\), the associated Newton flows tend to a degenerate flow. No common period or continuation of a particular attracting cycle is demanded by the supplied statement.

The conventional Newton-flow vector field for a fixed polynomial is
\[
 \frac{\dd z}{\dd t}=-\frac{f(z)}{f'(z)},
\]
for which the relaxed map is an explicit Euler step of length \(h\) away from singularities. For the proposed family the vector field also depends on \(h\). The source does not define ``degenerate flow'', the mode or domain of convergence, possible time rescaling, or the regularity of the parameter family. It also does not explicitly say that every \(f_h\) retains degree \(d\), how polynomial scaling is normalized, or whether multiple roots are allowed. These omissions prevent a complete limit theorem from being stated; no notion of degeneracy or convergence is substituted.

The registry labels this open/unverified and the summary reports no verified published resolution. The original problem is from 1992, and the supplied provenance records no comprehensive modern resolution check.
\subsection{Short English statement}
Can every polynomial with an unwanted attracting Newton cycle be deformed through polynomials whose relaxed Newton maps remain bad for every \(h>0\), ending in a degenerate Newton-flow limit as \(h\to0\)? The family and limit conditions, especially degeneracy, need their original definitions.
\subsection{Sources}
\begin{itemize}
\item Local registry: \url{data/unsolvedmath/problems.json}, permanent UnsolvedMath ID \(5300070\), source code catalogue entry; supplied statement, background, and research summary, updated 2026-08-06.
\item Scott Sutherland (1992), problem C1, in Stony Brook, \emph{Open Problems in Dynamical Systems}: \url{https://www.math.stonybrook.edu/open-problems-dynamical-systems}. Supplied source-TeX and PDF-text extraction accessed 2026-07-29; no comprehensive modern resolution check was available.
\item The summary refers to Sutherland's thesis-era Newton-map problems without providing a separate thesis or publication citation.
\end{itemize}
\end{document}
